% ПРИЛОЖЕНИЯ, без номер на раздел.
% По А.3: изходен текст на програми, формати на структури от данни, екранни форми,
% получени резултати и др. Изходният текст се ВМЪКВА с \lstinputlisting от src/,
% а не се преписва тук: преписаният код се разминава с истинския още при първата
% поправка. Всички листинги от diagrams/ са ПОРОДЕНИ от инструмента и се
% възпроизвеждат с  cmake --build build --target demo
\section*{ПРИЛОЖЕНИЯ}
\label{sec:appendix}
\addcontentsline{toc}{section}{ПРИЛОЖЕНИЯ}

\subsection*{Приложение А. Формат на вътрешния модел}
\addcontentsline{toc}{subsection}{Приложение А. Формат на вътрешния модел}

Вътрешният модел е записваем във външен вид, за да може да бъде сравняван и съхраняван
между отделните стъпки на пълния цикъл. Записът е JSON с два масива на най-горно ниво,
\code{classifiers} и \code{relations}. Класификаторът носи неквалифицирано и квалифицирано
име, вид, признак за абстрактност, шаблонни параметри, стереотипи, етикетирани стойности,
атрибути, операции, изброители и мястото в изходния текст, откъдето е дошъл. Атрибутът
носи име, тип както е изписан в кода, видимост, обхват, кратност и собствени етикетирани
стойности. Операцията носи име, тип на резултата, параметри, видимост, обхват и признаци за
виртуалност, чистота и константност. Връзката носи двата края, вид, етикет, кратност,
признак за виртуална база и видимост на базата.

Класификаторите се свързват по име, а не по указател. Изборът е продиктуван от изискването
моделът да може да бъде записан и прочетен обратно, а също и да бъде сравняван с друг модел
по стойност: имената преживяват запис във файл, указателите не.

\lstinputlisting[language=C++, caption={Дефиниция на вътрешния модел, \code{include/blueprint/model.hpp}}, label={lst:model}]{../include/blueprint/model.hpp}

\subsection*{Приложение Б. Примерен вход и получените от него диаграми}
\addcontentsline{toc}{subsection}{Приложение Б. Примерен вход и получените диаграми}

Приложението съдържа малката единица за превод, върху която е проведен пълният цикъл. Тя
обхваща наследяване, абстрактен базов клас, асоциация, агрегация, композиция и кратност
много, като всяка от тези конструкции се среща в нея точно по веднъж, за да сочи един
евентуален отказ към едно правило, а не към сплетени правила.

\lstinputlisting[language=C++, caption={Анотиран примерен вход, \code{examples/library.hpp}}, label={lst:library}]{../examples/library.hpp}

Следват двата текста, получени от този вход. И двата са породени от един и същ модел с
\code{blueprint diagram} и се различават само по нотацията.

\lstinputlisting[language={}, caption={Породена диаграма на класове в PlantUML}, label={lst:puml}]{../diagrams/library.class.puml}

\lstinputlisting[language={}, caption={Породена диаграма на класове в Mermaid}, label={lst:mmd}]{../diagrams/library.class.mmd}

Поведенческите модели се превеждат по същия начин. По-долу е диаграмата на случаите на
употреба в PlantUML и диаграмата на дейностите в Mermaid, тоест същите два модела, които
Фиг.~\ref{fig:usecase} и Фиг.~\ref{fig:activity} показват.

\lstinputlisting[language={}, caption={Породена диаграма на случаите на употреба в PlantUML}, label={lst:ucpuml}]{../diagrams/blueprint.usecase.puml}

\lstinputlisting[language={}, caption={Породена диаграма на дейностите в Mermaid}, label={lst:actmmd}]{../diagrams/roundtrip.activity.mmd}

\subsection*{Приложение В. Породена релационна схема и обратно породени заготовки}
\addcontentsline{toc}{subsection}{Приложение В. Породена релационна схема}

Схемата по-долу е получена от модела на Приложение Б при стратегия \term{една таблица за
йерархията}. Класовете \code{Book} и \code{Journal} са свити в таблицата \code{items},
която получава колона разграничител, а колоните им допускат празна стойност. Асоциацията с
кратност много между \code{Catalogue} и \code{LibraryItem} поражда свързваща таблица.
Името на класа пътува в коментар пред всяка таблица, а името на атрибута в коментар след
колоната, чието име се различава от неговото, за да не се налага при обратния преход да се
отгатва единствено число на английска дума.

\lstinputlisting[language=SQL, caption={Породена схема при стратегия „една таблица за йерархията“}, label={lst:sql}]{../diagrams/library.single-table.sql}

Обратният преход от схема към код затваря цикъла. Породеният заглавен файл съдържа само
декларации и носи макросите за анотиране, така че повторното извличане от него дава
съпоставим модел. Че той е действителен C++, се проверява в демонстрацията с
\code{clang -fsyntax-only}, а не се твърди.

\lstinputlisting[language=C++, caption={Заготовки на класове, породени обратно от схемата}, label={lst:skeleton}]{../diagrams/library_generated.hpp}

\subsection*{Приложение Г. Ключови места от изходния текст}
\addcontentsline{toc}{subsection}{Приложение Г. Ключови места от изходния текст}

Пълният изходен текст е в хранилището. Тук се вмъкват двете места, на които се взема
същинското решение, а не се пренася информация.

Първото е правилото, което превръща изписването на един член във вид на връзка. То е
съзнателно събрано в една функция, за да може да бъде прочетено наведнъж и оспорено
наведнъж.

\lstinputlisting[language=C++, firstline=133, lastline=212, firstnumber=133, caption={Правилото за собствеността, \code{src/reader.cpp}}, label={lst:ownership}]{../src/reader.cpp}

Второто е броенето на запазените елементи, което е числителят на дела от
Табл.~\ref{tab:roundtrip}. То се брои поелементно, а не се извежда от броя на разликите,
и коментарът в кода казва защо.

\lstinputlisting[language=C++, firstline=116, lastline=142, firstnumber=116, caption={Броене на запазените елементи, \code{src/diff.cpp}}, label={lst:preserved}]{../src/diff.cpp}
